\documentclass[twocolumn,11pt]{article}
%\documentclass[11pt]{article}
\usepackage[margin=1in]{geometry}
\usepackage{chemfig}
\usepackage{mhchem}
\usepackage{gensymb} % For \degree at least.

\title{Dean's Chemistry Chapter 8 \LaTeX\ in TeXstudio}
\author{Dean E. Peterson}
\date{April 16, 2026}
\renewcommand{\labelenumi}{(\alph{enumi})}

\begin{document}
\maketitle

\section{Chapter 8}
\subsection{Chapter 8 Examples}
\subsubsection{Example 8.2 - Solution}

Ammonium sulfate is important as a fertilizer. What is the hybridization of the sulfur atom in the sulfate ion, \ce{SO4^{2-}}?

32 valence electrons total, 24 after adding 4 single bonds, 0 after giving each O 3 lone pairs. \\[1ex]

$\left[\enspace
\chemfig{
    \charge{[circle]-60:6pt=\scriptsize{+2}}{S}
    ( -[:0]\charge{[circle]-60:6pt=\scriptsize{-1},0=\:,90=\:,-90=\:}{O} )
    ( -[:90]\charge{[circle]-60:6pt=\scriptsize{-1},0=\:,90=\:,180=\:}{O} )
    ( -[:180]\charge{[circle]-60:6pt=\scriptsize{-1},180=\:,90=\:,-90=\:}{O} )
    ( -[:270]\charge{[circle]-60:6pt=\scriptsize{-1},0=\:,180=\:,-90=\:}{O} )
}
\enspace\right]^{2-}$
, then maybe \\[5ex]
$\left[\enspace
\chemfig{
    \charge{[circle]-60:6pt=\scriptsize{0}}{S}
    ( =[:0]\charge{[circle]-60:6pt=\scriptsize{0},90=\:,-90=\:}{O} )
    ( -[:90]\charge{[circle]-60:6pt=\scriptsize{-1},0=\:,90=\:,180=\:}{O} )
    ( =[:180]\charge{[circle]-60:6pt=\scriptsize{0},90=\:,-90=\:}{O} )
    ( -[:270]\charge{[circle]-60:6pt=\scriptsize{-1},0=\:,180=\:,-90=\:}{O} )
}
\enspace\right]^{2-}$


\subsubsection{Example 8.3 - Solution}
Urea, $\ce{NH2C(O)NH2}$, is sometimes used as a source of nitrogen in fertilizers. What is the hybridization of the carbon atom in urea?

There are 24 valence electrons. 10 left after single bonds \\[2ex]

\chemfig{
    \charge{[circle]-60:6pt=\scriptsize{0},180=\:}{N}
    ( -[:90]H)
    ( -[:-90]H)
    -\charge{[circle]-60:6pt=\scriptsize{0}}{C}
    ( =[:90]\charge{[circle]-60:6pt=\scriptsize{0},0=\:,180=\:}{O} )
    -\charge{[circle]-60:6pt=\scriptsize{0},0=\:}{N}
    ( -[:90]H)
    ( -[:-90]H)
}

\subsubsection{Example 8.4 - Check Your Learning}
The long and the short of it is I want the Lewis structure(s) for \ce{NO2}.

\ce{NO2} has 17 total valence electrons, 13 after single bonds, 1 after giving 3 lone pairs to each oxygen. \\[1ex]

\chemfig{
    \charge{[circle]-70:6pt=\scriptsize{+2},90=\.}{N}
    ( -[:-150]\charge{[circle]-60:6pt=\scriptsize{-1},180=\:,-90=\:,90=\:}{O} )
    ( -[:-30]\charge{[circle]-60:6pt=\scriptsize{-1},90=\:,-90=\:,0=\:}{O} )
}
\\[4ex]
So...\\

\chemfig{
    \charge{[circle]-70:6pt=\scriptsize{+1},90=\.}{N}
    ( =[:-150]\charge{[circle]-60:6pt=\scriptsize{0},180=\:,-90=\:}{O} )
    ( -[:-30]\charge{[circle]-60:6pt=\scriptsize{-1},90=\:,-90=\:,0=\:}{O} )
}
\quad$\longleftrightarrow$\quad
\chemfig{
    \charge{[circle]-70:6pt=\scriptsize{+1},90=\.}{N}
    ( -[:-150]\charge{[circle]-60:6pt=\scriptsize{-1},180=\:,-90=\:,90=\:}{O} )
    ( =[:-30]\charge{[circle]-60:6pt=\scriptsize{0},-90=\:,0=\:}{O} )
}

\subsubsection{Exercise 8.15}

\begin{enumerate}

\item circular $\ce{S8}$ molecule, 48 valence, 32 after single bonds, or 4 per S with single bonds \\[2ex]
\chemfig{
    \charge{[circle]180=\:,-105=\:}{S}*8 (
    - \charge{[circle]-90=\:,-165=\:}{S}
    - \charge{[circle]-90=\:,-15=\:}{S}
    - \charge{[circle]0=\:,-75=\:}{S}
    - \charge{[circle]0=\:,75=\:}{S}
    - \charge{[circle]90=\:,15=\:}{S}
    - \charge{[circle]90=\:,165=\:}{S}
    - \charge{[circle]180=\:,105=\:}{S}
    -
    )
}


\item $\ce{SO2}$ molecule, 18 valence, 14 after single bonds, 2 after give each O 6.\\[2ex]

Copied from exercise 7.45\\[2ex]
\chemfig{
    \charge{[circle]-60:6pt=\scriptsize{-1},180=\:,90=\:,-90=\:}{O}
    -\charge{[circle]-60:6pt=\scriptsize{+1},-90=\:}{S}
    =\charge{[circle]-60:6pt=\scriptsize{0},90=\:,-90=\:}{O}
}
\quad $\longleftrightarrow$ \quad
\chemfig{
    \charge{[circle]-60:6pt=\scriptsize{0},90=\:,-90=\:}{O}
    =\charge{[circle]-60:6pt=\scriptsize{+1},-90=\:}{S}
    -\charge{[circle]-60:6pt=\scriptsize{-1},0=\:,90=\:,-90=\:}{O}
}\\[1ex]

\item $\ce{SO3}$ molecule, 24 valence, 18 after single bonds, 0 after give each O 6.\\[2ex]
\chemfig{
    \charge{[circle]-60:6pt=\scriptsize{-1},180=\:,90=\:,-90=\:}{O}
    - \charge{[circle]-60:6pt=\scriptsize{+2}}{S}
    ( =[:90] \charge{[circle]-60:6pt=\scriptsize{0},180=\:,0=\:}{O})
    - \charge{[circle]-60:6pt=\scriptsize{-1},0=\:,90=\:,-90=\:}{O}
}
\quad $\longleftrightarrow$ \quad \text{etc.}
\\[3ex]
\chemfig{
    \charge{[circle]-60:6pt=\scriptsize{0},90=\:,-90=\:}{O}
    = \charge{[circle]-60:6pt=\scriptsize{0}}{S}
    ( =[:90] \charge{[circle]-60:6pt=\scriptsize{0},180=\:,0=\:}{O})
    = \charge{[circle]-60:6pt=\scriptsize{0},90=\:,-90=\:}{O}
}
\\[4ex]
\chemfig{
    \charge{[circle]-60:6pt=\scriptsize{0},180=\:,90=\.,-90=\.}{O}
    = \charge{[circle]-60:6pt=\scriptsize{0}}{S}
    ( =[:90] \charge{[circle]-60:6pt=\scriptsize{0},180=\.,90=\:,0=\.}{O})
    = \charge{[circle]-60:6pt=\scriptsize{0},0=\:,90=\.,-90=\.}{O}
}\\[1ex]

\item $\ce{H2SO4}$ molecule (the hydrogen atoms are bonded to oxygen atoms),
32 valence, 20 after single bonds.\\[2ex]
\chemfig{
    \charge{[circle]-60:6pt=\scriptsize{+2}}{S}
    (-[:180]\charge{[circle]-90=\:,180=\:,90=\:,-60:6pt=\scriptsize{-1}}{O})
    (-[:0]\charge{[circle]-90=\:,0=\:,90=\:,-60:6pt=\scriptsize{-1}}{O})
    (-[:90]\charge{[circle]90=\:,180=\:,-60:6pt=\scriptsize{0}}{O}
    -[:0]\charge{[circle]-60:6pt=\scriptsize{0}}{H})
    (-[:-90]\charge{[circle]-90=\:,180=\:,-60:6pt=\scriptsize{0}}{O}
    -[:0]\charge{[circle]-60:6pt=\scriptsize{0}}{H})
}
\qquad \qquad
\chemfig{
    \charge{[circle]180=\:,-60:6pt=\scriptsize{+1}}{S}
    (-[:0]\charge{[circle]-90=\:,90=\:,-60:6pt=\scriptsize{0}}{O}
    -[:0]\charge{[circle]-90=\:,0=\:,90=\:,-60:6pt=\scriptsize{-1}}{O})
    (-[:90]\charge{[circle]90=\:,180=\:,-60:6pt=\scriptsize{0}}{O}
    -[:0]\charge{[circle]-60:6pt=\scriptsize{0}}{H})
    (-[:-90]\charge{[circle]-90=\:,180=\:,-60:6pt=\scriptsize{0}}{O}
    -[:0]\charge{[circle]-60:6pt=\scriptsize{0}}{H})
}
\\[4ex]
Book answer to exercise 7.63 was like:\\[3ex]
\chemfig{
    \charge{[circle]-60:6pt=\scriptsize{0}}{S}
    (=[:180]\charge{[circle]-90=\.,180=\:,90=\.,-60:6pt=\scriptsize{0}}{O})
    (=[:0]\charge{[circle]-90=\.,0=\:,90=\.,-60:6pt=\scriptsize{0}}{O})
    (-[:90]\charge{[circle]90=\:,180=\:,-60:6pt=\scriptsize{0}}{O}
    -[:0]\charge{[circle]-60:6pt=\scriptsize{0}}{H})
    (-[:-90]\charge{[circle]-90=\:,180=\:,-60:6pt=\scriptsize{0}}{O}
    -[:0]\charge{[circle]-60:6pt=\scriptsize{0}}{H})
}


\end{enumerate}



\end{document}
